\section{Background and Significance}

\subsection{Background}

In the intersection of cognitive science and artificial intelligence, numerical processing ability has long been regarded as a core indicator of intelligence. In recent years, Large Language Models (LLMs) based on the Transformer architecture—such as GPT-4 and Llama 3—have made breakthrough progress in natural language processing, logical reasoning, and code generation. However, despite their astounding general capabilities, these models frequently exhibit rudimentary errors in basic numerical comparison and arithmetic operations that are inconsistent with high-level intelligence. A typical case is the "9.11 vs. 9.9 comparison paradox," where leading frontier models often erroneously assert that "9.11 is larger" when asked "Which is larger, 9.11 or 9.9?"\cite{eval16x2025}.

This phenomenon is not isolated, nor is it a simple calculation error; rather, it represents a systematic cognitive bias. Interestingly, this bias bears a striking resemblance to the "Number Bias" recorded in human developmental psychology. During cognitive development, children and even some adults tend to erroneously apply the rules of natural numbers (integers) when processing rational numbers (fractions and decimals). This phenomenon is known as the Natural Number Bias (NNB) or Whole Number Bias (WNB)\cite{alibali2015}. For example, children often misjudge "9.11 > 9.9" because "11 > 9," or believe that a number with more decimal places must be larger.

Although academia has conducted extensive independent research on human cognitive deficits and neural network hallucinations, there is a lack of comparative research placing both within a unified framework. As systems based on statistical probability to predict the next token, do the internal mechanisms of LLMs reenact the dual-process mechanism of the human brain when processing numbers? Do the model's errors stem from physical defects in the tokenization mechanism\cite{singh2024}, or from schema interference in training data where non-mathematical contexts (such as dates or software versions) dominate? This project intends to compare the behavioral patterns of humans and LLMs in numerical processing tasks to explore the similarities, differences, and underlying mechanisms of their numerical biases.

\subsection{Significance}

This study possesses both significant theoretical value and application prospects.

Theoretical Level: This study helps deepen the understanding of the essence of intelligence. By treating LLMs as a "mirror" or "high-throughput agent" of human cognition, we can verify the applicability of classical theories in developmental psychology (such as Conceptual Change Theory and Dual Process Theory) within artificial neural networks. If it is confirmed that LLMs exhibit cognitive biases isomorphic to humans, this will provide new evidence for whether "connectionist" AI possesses human-like cognitive structures, while also offering a new, scalable simulation platform for cognitive science research\cite{vosniadou1994}.

Application Level: This study is crucial for enhancing the reliability and safety of LLMs. Currently, LLMs are accelerating their integration into fields requiring high numerical precision, such as financial analysis, educational tutoring, and scientific computing. The instability of models in basic numerical understanding constitutes a potential risk. By parsing the specific pathways through which models generate numerical bias—whether from the physical cutting of tokenizers or the overwhelming interference of "version number schemas" over "mathematical schemas" in semantic space—we can design more targeted intervention measures. This will not only guide the cleaning and ratio optimization of future training data but also facilitate the development of more effective alignment algorithms or prompt engineering strategies, thereby fundamentally mitigating logical hallucinations\cite{turpin2023} and promoting the development of AI towards greater rigor and explainability.
